\documentclass[12pt]{article}
\usepackage{sty/Packages}
\usepackage{sty/Presets}
\usepackage{sty/Commands}
\usepackage{sty/Misc}

\begin{document}

\begin{titlepage}
\begin{center} 
 {\large {\fontfamily{qcs}\textsc{Vietnam National University - Ho Chi Minh city}} \medskip \\
  {\large \fontfamily{qcs}\textsc{University of Technology}} \medskip \\
  \large {\fontfamily{qcs}\textsc{Faculty of Computer Science and Engineering}} \\
  \LARGE\bd{------------------------}

  \vspace{20pt}
  
  \includegraphics[scale = 0.3]{img/logo.png}
 }

\vspace{15pt}

 \bd{\LARGE\bd{------------------------------------------------------------} \\ 
      Parallel Computing (CO3068) \\
      Lab 4 Report \\ \medskip
      \Huge Distributed Simulation with Open MPI
      \LARGE \bd{------------------------------------------------------------} \\ 
    }

\vspace{15pt}

\large
\begin{otherlanguage}{vietnamese}
\begin{tabular}{rl}
    \bd{Instructor:} & Nguyễn Mạnh Thìn \\
    \bd{Student name:} & Nguyễn Văn An \\
    \bd{Student ID:} & 2352013
\end{tabular}
\end{otherlanguage}

\bigskip\bigskip\bigskip
\normalsize
\vfill
\textit{Ho Chi Minh city, \monthname, \the\year}
\end{center}
\end{titlepage}
\tableofcontents
\newpage

% --- SECTION 1: INTRODUCTION ---
\section{Introduction}

This lab focuses on the implementation of a distributed parallel simulation using the \textbf{Message Passing Interface (MPI)}. Building upon the previous exercises (Heat Diffusion, Radioactive Contamination, and Shock Wave propagation), the goal is to scale the simulation across multiple physical nodes in a distributed memory environment.

The primary objective is to implement two distinct communication strategies: \textbf{Synchronous} and \textbf{Asynchronous}, and analyze their performance impact. Specifically, we aim to demonstrate "Latency Hiding" by overlapping communication with computation to improve the overall runtime efficiency of the system.

The simulation operates on a $4000 \times 4000$ grid over 100 time steps. The workload is distributed using a row-wise domain decomposition strategy, where each MPI process is responsible for a horizontal "slab" of the grid and maintains "halo" (ghost) rows to exchange boundary data with neighbors.

% --- SECTION 2: IMPLEMENTATION STRATEGIES ---
\section{Implementation Strategies}

\subsection{Synchronous Implementation (Blocking)}
The synchronous approach utilizes standard blocking communication. At the end of each time step, every process pauses its computation to exchange halo rows with its neighbors.

\begin{itemize}
    \item \textbf{Mechanism:} Uses \texttt{MPI\_Sendrecv} to perform point-to-point communication.
    \item \textbf{Workflow:}
    \begin{enumerate}
        \item Compute all local rows (Interior and Boundaries).
        \item \textbf{Block} and wait for data exchange (Send to/Recv from neighbors).
        \item Proceed to the next iteration.
    \end{enumerate}
    \item \textbf{Disadvantage:} The CPU remains idle during the data transfer phase, leading to zero overlap between computation and communication.
\end{itemize}

\subsection{Asynchronous Implementation (Latency Hiding)}
The asynchronous approach leverages non-blocking MPI primitives to overlap the communication overhead with useful computational work. This is the core difference compared to the synchronous version.

\begin{itemize}
    \item \textbf{Mechanism:} Uses \texttt{MPI\_Isend} and \texttt{MPI\_Irecv} combined with \texttt{MPI\_Waitall}.
    \item \textbf{Latency Hiding Strategy:} To achieve true overlap, the computation kernel was split into two distinct phases:
    \begin{enumerate}
        \item \textbf{Phase A (Interior):} Rows that do not depend on neighbor data (Rows $2$ to $N-1$). These are safe to calculate immediately.
        \item \textbf{Phase B (Boundaries):} Rows that depend on halo data (Row $1$ and Row $N$). These must wait for communication to complete.
    \end{enumerate}
    \item \textbf{Workflow:}
    \begin{enumerate}
        \item Initiate non-blocking transfers (\texttt{MPI\_Isend} / \texttt{MPI\_Irecv}).
        \item \textbf{Compute Interior Rows} immediately. This computation happens \textit{while} the data is traveling over the network, effectively hiding the transfer latency.
        \item Call \texttt{MPI\_Waitall} to synchronize and ensure halos have arrived.
        \item Compute Boundary Rows.
    \end{enumerate}
\end{itemize}

% --- SECTION 3: KEY CODE IMPLEMENTATION ---
\section{Key Code Implementation}

This section highlights the critical differences in the main simulation loop between the Synchronous and Asynchronous versions.

\subsection{Synchronous Logic}
In the synchronous version, the physics computation and halo exchange happen sequentially. The process must finish all math before it can start talking to neighbors.

\begin{lstlisting}[style=Cpp, caption={Synchronous Halo Exchange using MPI\_Sendrecv}]
// 1. Heat Diffusion
compute_heat(my_rows, heat_curr, heat_next);

// 2. Blocking Exchange (CPU waits here)
exchange_halos(heat_next, my_rows, up_proc, down_proc);

// 3. Swap Pointers
std::swap(heat_curr, heat_next);
\end{lstlisting}

\subsection{Asynchronous Logic (Optimized)}
In the optimized asynchronous version, the computation is split. The \texttt{compute\_heat\_interior} function processes the majority of the grid while the MPI requests are active in the background.

\begin{lstlisting}[style=Cpp, caption={Asynchronous Latency Hiding Pattern}]
// A. Start Communication (Returns immediately)
start_exchange(heat_curr, my_rows, up_proc, down_proc, reqs); 

// B. Compute INTERIOR (Independent of Halos)
// [OVERLAP: CPU works here while Data moves over Network]
compute_heat_interior(my_rows, heat_curr, heat_next);         

// C. Wait for Halos (Synchronization Point)
MPI_Waitall(4, reqs, MPI_STATUSES_IGNORE);                    

// D. Compute BOUNDARIES (Now that we have valid halos)
compute_heat_boundaries(my_rows, heat_curr, heat_next);       

// E. Swap
std::swap(heat_curr, heat_next);
\end{lstlisting}

The \texttt{start\_exchange} function utilizes non-blocking primitives:
\begin{lstlisting}[style=Cpp, caption={Non-blocking MPI Calls}]
void start_exchange(float* buf, size_t rows, int up, int down, MPI_Request* reqs) {
    int tag = 0;
    // Initiate Receives (Irecv)
    MPI_Irecv(&buf[idx(0, 0)], paddedCsvSize, MPI_FLOAT, up, tag, MPI_COMM_WORLD, &reqs[0]);
    MPI_Irecv(&buf[idx(pad_sz + rows, 0)], ..., down, tag, MPI_COMM_WORLD, &reqs[1]);

    // Initiate Sends (Isend)
    MPI_Isend(&buf[idx(pad_sz, 0)], paddedCsvSize, MPI_FLOAT, up, tag, MPI_COMM_WORLD, &reqs[2]);
    MPI_Isend(&buf[idx(pad_sz + rows - 1, 0)], ..., down, tag, MPI_COMM_WORLD, &reqs[3]);
}
\end{lstlisting}

% --- SECTION 4: EXPERIMENTAL SETUP ---
\section{Experimental Setup}

The experiments were conducted on a heterogeneous distributed cluster consisting of 2 physical nodes connected via a standard Local Area Network (LAN).

\subsection{Hardware Configuration}
The cluster is composed of two distinct machines:
\begin{itemize}
    \item \textbf{Node 1 (Master/Rank 0):}
    \begin{itemize}
        \item CPU: 6 Cores
        \item RAM: 16 GB
    \end{itemize}
    \item \textbf{Node 2 (Worker):}
    \begin{itemize}
        \item CPU: 4 Cores
        \item RAM: 4 GB
    \end{itemize}
\end{itemize}

\subsection{Software Environment}
\begin{itemize}
    \item \textbf{Library:} Open MPI 5.0.8
    \item \textbf{Compilation:} Both implementations were compiled with the \texttt{-O2} optimization flag and native architecture tuning. This flag was chosen to ensure a balanced ratio between computation and communication, as aggressive optimization (e.g., \texttt{-O3}) can reduce computation time so significantly that communication overhead dominates (Amdahl's Law), making overlap difficult to observe.
    \begin{itemize}
        \item \texttt{mpicxx -O2 -march=native main.cpp -o main}
        \item \texttt{mpicxx -O2 -march=native main\_async.cpp -o main\_async}
    \end{itemize}
\end{itemize}

% --- SECTION 5: RESULTS ---
\section{Results}

The simulation was executed with \textbf{4 MPI Processes} distributed across the two nodes described above. The total execution time was measured for 3 independent runs for both implementations.

\begin{table}[H]
    \centering
    \caption{Runtime Performance Comparison (Seconds)}
    \begin{tabular}{lcc}
        \hline
        \textbf{Run \#} & \textbf{Synchronous (Blocking)} & \textbf{Asynchronous (Optimized)} \\
        \hline
        1 & 2.32167 s & 1.60134 s \\
        2 & 2.32432 s & 1.54456 s \\
        3 & 2.30868 s & 1.57610 s \\
        \hline
        \textbf{Average} & \textbf{2.318 s} & \textbf{1.574 s} \\
        \hline
    \end{tabular}
\end{table}

% --- SECTION 6: ANALYSIS ---
\section{Analysis and Discussion}

\subsection{Performance Comparison}
The Asynchronous implementation achieved an average runtime of \textbf{1.574 seconds}, compared to \textbf{2.318 seconds} for the Synchronous version.
\begin{itemize}
    \item \textbf{Speedup:} $2.318 / 1.574 \approx \mathbf{1.47x}$
    \item \textbf{Improvement:} The asynchronous version is approximately \textbf{32\% faster}.
\end{itemize}

\subsection{Latency Hiding Effectiveness}
The significant speedup confirms that the \textbf{Latency Hiding} strategy was highly effective in this distributed environment.
\begin{itemize}
    \item In the \textbf{Synchronous} model, the total time per iteration is largely additive: $T_{total} = T_{compute} + T_{comm}$. The network latency is fully exposed, as the CPU must idle while waiting for the network card to transmit halo data.
    \item In the \textbf{Asynchronous} model, the CPU processed the "Interior" rows which constitute the vast majority of the grid, \textit{while} the MPI data was in flight. Since the time required to compute the interior ($T_{interior}$) exceeded the communication time ($T_{comm}$), the network cost was effectively completely hidden. The effective iteration time approached $T_{compute}$.
\end{itemize}

\subsection{Impact of Optimization Flags}
The choice of the \texttt{-O2} flag played a critical role in these results. In preliminary tests using \texttt{-O3}, the compiler optimized the stencil computation loop so aggressively (using AVX vectorization) that $T_{compute}$ became negligible. When $T_{compute}$ is extremely small, it cannot effectively "cover" the cost of $T_{comm}$, leading to reduced scalability. By using \texttt{-O2}, we maintained a computation load sufficient to demonstrate the architectural benefits of the asynchronous design.

% --- SECTION 7: CONCLUSION ---
\section{Conclusion}

This lab successfully demonstrated the deployment of a hybrid physics simulation on a distributed cluster. By refactoring the standard stencil algorithm to separate independent (interior) and dependent (boundary) data regions, we were able to utilize \texttt{MPI\_Isend} and \texttt{MPI\_Irecv} to effectively hide communication latency.

The resulting \textbf{1.47x speedup} validates that overlapping communication with computation is a critical optimization technique for high-performance distributed computing, particularly in environments where network latency is a potential bottleneck.

\end{document}